\chapter{Calculer et justifier les dérivées}
\section*{Exercice 9.1 -- Somme et puissances}
Dériver $f(x)=3x^4-2x^{-2}+7$ en précisant son domaine.

\section*{Exercice 9.2 -- Produit}
Dériver $(x^2+1)\sqrt{x+2}$ sur l'intérieur de son domaine.

\section*{Exercice 9.3 -- Quotient}
Dériver $f(x)=(2x+1)/(x-3)$, puis retrouver le résultat en réécrivant $f$.

\section*{Exercice 9.4 -- Composition affine}
Dériver $g(x)=(3x-2)^5$ et $h(x)=1/(2x+1)$.

\section*{Exercice 9.5 -- Démontrer une règle}
À partir de $(fg)(a+h)-(fg)(a)$, retrouver la règle de dérivation d'un produit.

\section*{Exercice 9.6 -- Une erreur à réfuter}
Un élève affirme $(f/g)'=f'/g'$. Donner un contre-exemple où les deux quotients sont définis.
